\documentclass[10pt,a4paper]{amsart}
\usepackage[margin=19mm]{geometry}
\usepackage{amsmath,amssymb,mathtools}
\usepackage[scr=rsfs,cal=euler]{mathalfa}
\usepackage{xfrac}
\usepackage[unicode,hidelinks,bookmarks=false]{hyperref}
\newcommand{\ie}{i.e.\ }
\newcommand{\cf}{cf.\ }
\newcommand{\resp}{respectively\ }
\newcommand{\Subsection}{\subsection}
\providecommand{\nobreakdash}{\nobreak\hbox{-}}
\setlength{\emergencystretch}{2em}
\multlinegap0pt
\usepackage{bm}
\usepackage{bbm}
\usepackage{dsfont}
\usepackage{xspace}
\usepackage{cancel}
\usepackage{mathtools}
\usepackage{amsthm}
\makeatletter
\@ifundefined{definition}{\newtheorem{definition}{Definition}}{}
\@ifundefined{lemma}{\newtheorem{lemma}[definition]{Lemma}}{}
\makeatother
\title[Lorentz meets Ptolemy]{Lorentz meets Ptolemy}
\author{Felix Rott, Zhe-Feng Xu, Matteo Zanardini}
\date{Source: arXiv:2601.21596v1 (CC BY 4.0)}
\begin{document}
\maketitle

\noindent\emph{Source and licence.} Lorentz meets Ptolemy. Version arXiv:2601.21596v1 (CC BY 4.0), distributed under the Creative Commons Attribution 4.0 International licence, as recorded in the arXiv licence metadata for this exact version and shown on the abstract page https://arxiv.org/abs/2601.21596v1. Full rights notice: licenses/arxiv-2601.21596-CC-BY-4.0.txt.\par\smallskip

\noindent\textit{Editorial scope.} Counted unit: the Lemma whose source label is \texttt{lem: Ptolemy imply maximising} in the article (source0.tex lines 1147--1151, source label \texttt{lem: Ptolemy imply maximising}), delivered together with the article's own complete original proof of it (source0.tex lines 1153--1176, one \texttt{proof} environment of 24 lines). No mathematics is written, repaired or re-derived: statement and proof are the article's own text line for line. Required context retained uncounted: none. Every definition, display and label of those lines is retained unchanged. Omitted from this package and not redistributed: 1--1146, 1152--1152, 1177--1360. Nothing inside a retained excerpt is omitted or altered: every retained body is delivered exactly as the article's lines are written, so no omission receipt of any kind accompanies this package. Citations: the retained excerpts carry no citation command at all, so no bibliography entry of the article is transcribed and no reference is redistributed. Reference adaptation: none was needed and none was made; every reference inside the delivered unit resolves inside it, so no unresolved reference or citation is produced. Printed slips kept literal: no printed word, symbol, hypothesis or number of the counted statement or of its proof was altered. Layout and class: the article's own class is \texttt{article} with packages including inputenc, geometry, amssymb, amsmath, bm, mathtools, amsthm, babel, lmodern, microtype, euscript, hyperref, xcolor, todonotes; the wrapper is the lane's pinned \texttt{amsart} editorial wrapper at 10pt on A4, which restates the theorem-like environments the retained text needs and adds the article's own macro definitions that the retained text needs (none), with extra wrapper packages (bm, bbm, dsfont, xspace, cancel, mathtools). The delivered numbering is the unit's own. Assets: no retained span contains a figure environment, a graphics-inclusion command, a PDF-page inclusion, a table or a TikZ picture; no image file is copied into this package, the package declares no asset dependency and no image file is redistributed with it. Licence: version arXiv:2601.21596v1 of this article is distributed under the Creative Commons Attribution 4.0 International licence, as recorded in the arXiv licence metadata for this exact version and shown on the abstract page https://arxiv.org/abs/2601.21596v1. A full rights notice is delivered at licenses/arxiv-2601.21596-CC-BY-4.0.txt. Characters: every retained excerpt is pure ASCII, so the delivered engine, XeLaTeX with its default Latin Modern text font, reports no missing character.
\begin{lemma}[Globally maximising geodesics and empty cut locus]
\label{lem: Ptolemy imply maximising}
Let $(M,g)$ be a Ptolemaic spacetime, then all timelike geodesics are maximising geodesics. 
In particular, the timelike cut locus of any point is empty and there exists at most one future-directed maximising timelike geodesic connecting any two points.
\end{lemma}

\begin{proof}
We concentrate on future-directed geodesics since for past-directed geodesics the argument is the same. 
Let $\gamma :[0,b) \rightarrow M$ be a future-directed timelike geodesic parametrized by proper time and suppose that $0 \leq a < b$ is the parameter value of the first point along $\gamma$ in the (future) timelike cut locus of $\gamma(0) =p$. 
It is clear that $0<a$ and $\ell(\gamma_0, \gamma_t) = t$ for any $0 \leq t \leq a$, but 
$\ell(\gamma_0, \gamma_c) <c$ for all $a<c < b$. Let $\varepsilon>0$ be sufficiently small such that $a+ \varepsilon <b$ and $\gamma : [a- \varepsilon, a + \varepsilon ] \rightarrow M$ is a maximising timelike geodesic parametrized by proper time between its endpoints. 
Applying the Ptolemaic inequality to $\gamma(0) \ll \gamma(a-\varepsilon) \ll \gamma(a) \ll \gamma(a+ \varepsilon)$, we get
\[
2a \varepsilon \geq (a-\varepsilon) \varepsilon + \varepsilon \, \ell( \gamma_0, \gamma _{a+ \varepsilon}) \, ,
\]
and so we get that $a+ \varepsilon \geq \ell(\gamma_0, \gamma_{a+ \varepsilon})$. A straightforward application of the reverse triangle inequality tells that actually $a+ \varepsilon = \ell(\gamma_0, \gamma_{a+ \varepsilon})$, which contradicts the fact that $\gamma_a$ is in the timelike cut locus of $\gamma_0$. 

Suppose by contradiction that there exist two future-directed maximising timelike geodesics $\sigma, \gamma$ connecting $x$ to $y \in I^+(x)$. 
Up to reparametrization we can suppose that $\sigma, \gamma : [0,1] \rightarrow M$. 
By the first part of the proof, the geodesic $\gamma$ can be extended to the interval $[0, 1 + \varepsilon]$ in such a way that it remains a future-directed maximising timelike geodesic for some $\varepsilon >0$ sufficiently small. 
Consider
\[
\eta _t : = \begin{cases}
\sigma_t, & t \in [0,1] \, ,\\
\gamma_t, & t \in [1, 1+ \varepsilon] \, .
\end{cases}
\]
It is clear that $\eta$ is a future-directed maximising timelike geodesic between $\gamma_0$ and $\gamma _{1+ \varepsilon}$ since $\ell(\gamma_0, \gamma_{1+\varepsilon})$ coincides with the Lorentzian length of $\eta$. 
In particular, $\eta$ is a smooth curve and so $\dot \gamma _1= \dot \sigma _1$. Hence $\gamma= \sigma$ follows by a continuity argument.
\end{proof}

\end{document}
